%% =====================================================================
%%  Courses 1--5: Python, Go, Rust, Scala, Java
%% =====================================================================

%% ============================ COURSE 1: PYTHON =======================
\course{1}{EPL-201}{Python Programming for AI \& Enterprise}{3--0--2--4}{Basic programming logic, computer fundamentals}
{From Pythonic fundamentals to production AI services, data pipelines and agents}

\csec{Core Competencies}
\begin{itemize}
\item Write clean, typed, tested and idiomatic Python following PEP 8 and modern packaging practice.
\item Process data at scale with NumPy, pandas and Polars, and use concurrency (asyncio, multiprocessing, free-threaded Python) correctly.
\item Build secure web APIs, AI services and LLM-agent applications with FastAPI, Pydantic and model SDKs.
\item Ship Python to production: packaging, performance tuning, containers, CI/CD and supply-chain security.
\end{itemize}

\csec{Five-Unit Curriculum}
\unit{I}{Python Foundations and Pythonic Thinking}{9}{%
Python 3 runtime and interpreters; variables, data types and operators; control flow; functions, arguments and scope;
strings, lists, tuples, sets and dictionaries; comprehensions; files and context managers; exceptions;
modules and packages; virtual environments and the uv/pip toolchain; Jupyter and VS Code workflows.}
{Builds a command-line expense tracker with file storage, input validation and error handling, packaged as a module.}
\unit{II}{Object-Oriented, Functional and Typed Python}{9}{%
Classes, inheritance, composition, dunder methods and dataclasses; protocols and abstract base classes;
iterators, generators and decorators; functional tools (map, filter, functools, itertools);
type hints, generics and static checking (mypy/pyright); unit testing with pytest, fixtures and mocking; code quality with Ruff.}
{Delivers a typed library-management package with dataclasses, decorators, 90\% pytest coverage and a clean Ruff/mypy report.}
\unit{III}{Data, Scientific Python and Concurrency}{9}{%
NumPy arrays and vectorisation; pandas and Polars for data wrangling; visualisation with Matplotlib and Plotly;
SQL access with SQLAlchemy and DuckDB; concurrency models: threads, the GIL and free-threaded Python, multiprocessing, asyncio;
profiling with cProfile and memory tools.}
{Produces a sales-analytics pipeline in Polars + DuckDB with charts, and an asyncio web scraper benchmarked against a threaded version.}
\unit{IV}{Web APIs and AI Engineering with Python}{9}{%
HTTP and REST; FastAPI routing, dependency injection and async endpoints; Pydantic data validation; authentication (OAuth2, JWT);
background tasks and queues (Celery, Redis); calling ML models and LLM APIs; structured outputs and tool calling;
building retrieval-augmented and agent applications; MCP servers in Python.}
{Ships a FastAPI service with JWT auth that serves an ML model and an LLM-powered assistant using tool calls and a Python MCP server.}
\unit{V}{Production Python: Performance, Packaging and Security}{9}{%
Project layout and pyproject.toml; packaging and publishing wheels; dependency locking; performance: vectorisation, Numba, Cython and
Rust extensions via PyO3; logging and observability; Docker images for Python; CI/CD with GitHub Actions;
security: secrets, dependency scanning (pip-audit), SBOMs and secure coding.}
{Publishes a versioned package to a private index with a Rust-accelerated hot path, Docker image, CI pipeline and a clean pip-audit report.}

\Needspace{16\baselineskip}\csec{Industry Tools \& Ecosystem}
\begin{toolstable}
Toolchain & Python 3.13/3.14, uv, pip, Poetry, Ruff, mypy, pyright, pytest, pre-commit \\
Data \& science & NumPy, pandas, Polars, DuckDB, SQLAlchemy, Matplotlib, Plotly, Jupyter \\
Web \& services & FastAPI, Pydantic, Django, Flask, Celery, Redis, Uvicorn \\
AI \& agents & scikit-learn, PyTorch, Hugging Face, Anthropic/OpenAI SDKs, LangGraph, MCP Python SDK \\
Performance & Numba, Cython, PyO3 / maturin, cProfile, py-spy \\
Delivery \& security & Docker, GitHub Actions, pip-audit, Bandit, CycloneDX SBOM \\
\end{toolstable}

\csec{Hands-on Laboratory (30 Hours)}
\labs
{CLI tools with files, exceptions and argparse}
{OOP design with dataclasses and protocols}
{Generators, decorators and typed utilities with pytest}
{Data wrangling with Polars and DuckDB}
{asyncio vs threads vs multiprocessing benchmark}
{FastAPI REST service with JWT and Pydantic}
{LLM tool-calling assistant and a Python MCP server}
{Packaging, Docker and CI with security scanning}

\Needspace{12\baselineskip}\csec{Course Outcomes}
\begin{cotable}
CO1 & Explain Python's data model, runtime and idiomatic constructs. & L2 Understand \\
CO2 & Apply object-oriented, functional and typed programming with automated tests. & L3 Apply \\
CO3 & Analyse data and choose the right concurrency model for a workload. & L4 Analyse \\
CO4 & Build secure web APIs and AI/LLM applications in Python. & L3 Apply \\
CO5 & Create a packaged, optimised, secure and deployable Python product. & L6 Create \\
\end{cotable}

\capstone{\textbf{AI-Powered Business Insights Service:} a FastAPI platform that ingests company data with Polars, serves forecasts from an ML model,
answers questions through an LLM agent with tool calls, and ships as a tested, containerised, security-scanned package with CI/CD.}

\csec{Aligned Careers}
{\raggedright
\role{Python Developer}\role{AI / ML Engineer}\role{Data Engineer}\role{Backend Engineer}
\role{Automation Engineer}\role{Agentic AI Developer}
\par}

\Needspace{12\baselineskip}\csec{Where Python Fits: 2026--2035 Outlook}
\begin{outlooktable}
AI, ML \& agents & First choice for model building, LLM applications, agents and evaluation; expected to remain the language of AI. \\
Data \& analytics & Core language for data pipelines, notebooks and analytics, alongside SQL. \\
Automation \& DevOps & Scripting, testing, infrastructure automation and AI-driven operations. \\
Performance trend & Free-threaded Python and Rust extensions (PyO3) close the speed gap for production workloads. \\
\end{outlooktable}

\Needspace{15\baselineskip}\csec{Recommended Industry Certifications}
\begin{certtable}
PCEP -- Certified Entry-Level Python Programmer & OpenEDG Python Institute & Entry & I \\
PCAP -- Certified Associate Python Programmer & OpenEDG Python Institute & Associate & I--III \\
PCPP1 / PCPP2 -- Certified Professional Python Programmer & OpenEDG Python Institute & Professional & II--V \\
Databricks Certified Associate Developer for Apache Spark (Python) & Databricks & Associate & III \\
NPTEL/SWAYAM: The Joy of Computing using Python & IIT faculty, Govt. of India & Academic & I, II \\
\end{certtable}

\csec{Text \& Reference Books}
\begin{refs}
\item L. Ramalho, \emph{Fluent Python}, 2nd ed., O'Reilly, 2022.
\item E. Matthes, \emph{Python Crash Course}, 3rd ed., No Starch Press, 2023.
\item B. Slatkin, \emph{Effective Python: 125 Specific Ways to Write Better Python}, 3rd ed., Addison-Wesley, 2024.
\item W. McKinney, \emph{Python for Data Analysis}, 3rd ed., O'Reilly, 2022.
\end{refs}

%% ============================ COURSE 2: GO ===========================
\course{2}{EPL-202}{Go Programming for Cloud-Native Systems}{3--0--2--4}{One programming language, networking basics}
{Simple, fast, concurrent services for cloud platforms, DevOps and AI infrastructure}

\csec{Core Competencies}
\begin{itemize}
\item Write idiomatic Go with clear types, interfaces, generics, explicit error handling and table-driven tests.
\item Design safe concurrent programs with goroutines, channels, select, context and the sync package.
\item Build high-performance HTTP and gRPC microservices with databases, observability and security.
\item Deliver cloud-native tools: CLIs, Kubernetes controllers, AI gateways and MCP servers.
\end{itemize}

\csec{Five-Unit Curriculum}
\unit{I}{Go Foundations}{9}{%
Go philosophy and toolchain (go build, run, fmt, vet); packages and modules; variables, constants and basic types;
control flow; functions, multiple returns and closures; arrays, slices and maps; strings, runes and UTF-8;
pointers; explicit error values and wrapping.}
{Builds a log-file analyser CLI that parses large files with slices and maps and reports errors with wrapping.}
\unit{II}{Types, Interfaces, Generics and Testing}{9}{%
Structs, methods and embedding; interfaces and implicit satisfaction; type switches; generics with type parameters and constraints;
package design and visibility; table-driven unit tests, benchmarks and fuzzing; examples and documentation;
dependency management with Go modules.}
{Delivers a generic, interface-based inventory library with table-driven tests, fuzz tests and benchmarks.}
\unit{III}{Concurrency in Go}{9}{%
Goroutines and the Go scheduler; channels (buffered, unbuffered) and select; worker pools, pipelines and fan-in/fan-out;
context for cancellation and deadlines; sync.Mutex, WaitGroup, Once and atomics; errgroup;
race detection; common concurrency bugs and how to avoid them.}
{Builds a concurrent web crawler with a worker pool, rate limiting and context cancellation that passes the race detector.}
\unit{IV}{Network Services and Microservices}{9}{%
net/http servers, routing and middleware; JSON APIs; gRPC and Protocol Buffers; database access with database/sql and pgx;
configuration and graceful shutdown; authentication; structured logging (slog); metrics and tracing with OpenTelemetry and Prometheus.}
{Ships a REST + gRPC order microservice on PostgreSQL with middleware, graceful shutdown and Prometheus metrics.}
\unit{V}{Cloud-Native Go and AI Infrastructure}{9}{%
Container-friendly Go builds (static binaries, distroless images); Kubernetes client-go and operators;
CLI tools with Cobra; profiling with pprof and performance tuning; security (govulncheck, secrets, TLS);
Go for AI infrastructure: LLM gateways, vector-search services and MCP servers.}
{Deploys a Kubernetes operator and an LLM API gateway in Go with caching, rate limits, pprof profiling and a clean govulncheck scan.}

\Needspace{16\baselineskip}\csec{Industry Tools \& Ecosystem}
\begin{toolstable}
Toolchain & Go 1.2x, go test, go vet, gofmt, golangci-lint, govulncheck, Delve \\
Web \& RPC & net/http, chi, Gin, Echo, gRPC-Go, Protocol Buffers, Connect \\
Data & database/sql, pgx, sqlc, Redis, NATS, Kafka (franz-go) \\
Cloud-native & Docker, Kubernetes client-go, Kubebuilder, Helm, Terraform \\
Observability & slog, OpenTelemetry, Prometheus, Grafana, pprof \\
AI infrastructure & MCP Go SDK, LLM provider Go SDKs, Weaviate / Milvus Go clients \\
\end{toolstable}

\csec{Hands-on Laboratory (30 Hours)}
\labs
{Go toolchain, modules and a CLI with Cobra}
{Slices, maps and error handling exercises}
{Interfaces and generics with table-driven tests}
{Goroutines, channels and worker pools}
{Context, cancellation and race detection}
{REST + gRPC microservice on PostgreSQL}
{Observability with slog, OpenTelemetry and pprof}
{Kubernetes operator and an LLM gateway}

\Needspace{12\baselineskip}\csec{Course Outcomes}
\begin{cotable}
CO1 & Explain Go's design, types, slices, maps and error model. & L2 Understand \\
CO2 & Apply interfaces, generics and testing to build reusable packages. & L3 Apply \\
CO3 & Analyse and build correct concurrent programs with goroutines and channels. & L4 Analyse \\
CO4 & Build observable HTTP and gRPC microservices with persistent storage. & L3 Apply \\
CO5 & Create cloud-native Go tools and AI infrastructure services for Kubernetes. & L6 Create \\
\end{cotable}

\capstone{\textbf{Cloud-Native AI Gateway Platform:} a Go platform that routes, caches and rate-limits requests to several LLM providers,
exposes REST and gRPC APIs, runs on Kubernetes with an operator for configuration, and provides full metrics and tracing.}

\csec{Aligned Careers}
{\raggedright
\role{Go Backend Engineer}\role{Cloud-Native Engineer}\role{Site Reliability Engineer}
\role{Platform / DevOps Engineer}\role{AI Infrastructure Engineer}
\par}

\Needspace{12\baselineskip}\csec{Where Go Fits: 2026--2035 Outlook}
\begin{outlooktable}
Cloud-native platforms & Kubernetes, Docker, Terraform and much of the cloud-native ecosystem are written in Go; expected to stay central. \\
Microservices \& APIs & Fast, low-memory services with simple deployment as single binaries. \\
AI infrastructure & Gateways, agent hosting, vector-search services and MCP servers that must be reliable at scale. \\
DevOps \& SRE & CLIs, automation and observability tooling. \\
\end{outlooktable}

\Needspace{15\baselineskip}\csec{Recommended Industry Certifications}
\begin{certtable}
Certified Kubernetes Application Developer (CKAD) & CNCF / Linux Foundation & Professional & V \\
Certified Kubernetes Administrator (CKA) & CNCF / Linux Foundation & Professional & V \\
Google Cloud Professional Cloud Developer & Google Cloud & Professional & IV, V \\
HashiCorp Certified: Terraform Associate & HashiCorp & Associate & V \\
GitHub portfolio assessment (no official Go language exam exists) & AlgoProfessor / institution & Project-based & I--V \\
\end{certtable}

\csec{Text \& Reference Books}
\begin{refs}
\item A. A. A. Donovan and B. W. Kernighan, \emph{The Go Programming Language}, Addison-Wesley, 2015.
\item J. Bodner, \emph{Learning Go}, 2nd ed., O'Reilly, 2024.
\item K. Cox-Buday, \emph{Concurrency in Go}, O'Reilly, 2017.
\item T. Harsanyi, \emph{100 Go Mistakes and How to Avoid Them}, Manning, 2022.
\end{refs}

%% ============================ COURSE 3: RUST =========================
\course{3}{EPL-203}{Rust Programming for Safe Systems}{3--0--2--4}{C or any compiled language, data structures}
{Memory-safe, high-performance software for systems, security, edge and AI infrastructure}

\csec{Core Competencies}
\begin{itemize}
\item Reason confidently about ownership, borrowing and lifetimes to write memory-safe code without a garbage collector.
\item Model problems with enums, traits, generics and pattern matching, with robust error handling.
\item Write fearless concurrent and async programs with threads, channels and Tokio.
\item Build production Rust: web services, WebAssembly, embedded firmware, Python extensions and AI inference.
\end{itemize}

\csec{Five-Unit Curriculum}
\unit{I}{Rust Foundations and Ownership}{9}{%
Why Rust: memory safety, performance and security; rustup, Cargo and crates.io; variables, mutability and data types;
functions and control flow; ownership, moves and copies; references and borrowing; slices; the stack and the heap;
reading compiler errors as guidance.}
{Builds a text-statistics CLI with Cargo that processes files using borrowed slices and passes Clippy with no warnings.}
\unit{II}{Types, Traits and Error Handling}{9}{%
Structs and methods; enums and pattern matching; Option and Result; error handling with ?, thiserror and anyhow;
traits, trait objects and generics; lifetimes in functions and structs; modules, crates and workspaces;
unit, integration and documentation tests.}
{Delivers a bank-ledger library modelled with enums and traits, typed errors and full unit and integration tests.}
\unit{III}{Collections, Iterators and Smart Pointers}{9}{%
Vec, String, HashMap and BTreeMap; iterators and closures; iterator adapters for zero-cost data processing;
Box, Rc, Arc, RefCell and interior mutability; Deref and Drop; macros (declarative); performance measurement with Criterion;
idiomatic API design.}
{Builds a fast CSV analytics engine with iterator chains and smart pointers, benchmarked with Criterion against a Python version.}
\unit{IV}{Fearless Concurrency, Async and Interoperability}{9}{%
Threads, message passing and shared state with Mutex and Arc; Send and Sync; Rayon for data parallelism;
async/await and the Tokio runtime; async I/O and networking; unsafe Rust and its rules;
FFI with C; Python extensions with PyO3 and maturin.}
{Ships a Tokio-based concurrent TCP chat server and a PyO3 extension that speeds up a Python hot loop.}
\unit{V}{Production Rust: Web, WebAssembly, Embedded and AI}{9}{%
Web services with Axum and SQLx; WebAssembly with wasm-bindgen; embedded Rust (embedded-hal, Embassy) on microcontrollers;
AI and ML in Rust (Candle, Burn, ONNX Runtime bindings); observability with tracing;
supply-chain security (cargo-audit, cargo-deny); Rust in the Linux kernel and safety-critical industries.}
{Deploys an Axum REST service with SQLx, a WebAssembly front-end module and an on-device inference demo, with a clean cargo-audit report.}

\Needspace{16\baselineskip}\csec{Industry Tools \& Ecosystem}
\begin{toolstable}
Toolchain & rustup, Cargo, Clippy, rustfmt, rust-analyzer, Miri, Criterion \\
Async \& web & Tokio, Axum, Actix Web, Hyper, Reqwest, SQLx, Serde \\
Parallel \& systems & Rayon, Crossbeam, nix, PyO3 / maturin, cbindgen \\
WebAssembly \& embedded & wasm-bindgen, wasm-pack, Leptos, embedded-hal, Embassy, probe-rs \\
AI \& data & Candle, Burn, ort (ONNX Runtime), Polars (Rust core), Hugging Face tokenizers \\
Security & cargo-audit, cargo-deny, cargo-fuzz, tracing \\
\end{toolstable}

\csec{Hands-on Laboratory (30 Hours)}
\labs
{Cargo projects and ownership exercises}
{Borrowing, slices and lifetimes puzzles}
{Enums, traits and typed error handling}
{Iterators, closures and Criterion benchmarks}
{Threads, Arc/Mutex and Rayon parallelism}
{Async TCP server with Tokio}
{PyO3 extension and C FFI}
{Axum + SQLx service, WebAssembly and embedded demo}

\Needspace{12\baselineskip}\csec{Course Outcomes}
\begin{cotable}
CO1 & Explain ownership, borrowing and lifetimes and why they guarantee memory safety. & L2 Understand \\
CO2 & Apply enums, traits, generics and error handling to model real problems. & L3 Apply \\
CO3 & Analyse performance and write efficient code with iterators and smart pointers. & L4 Analyse \\
CO4 & Build safe concurrent, async and cross-language (FFI, PyO3) components. & L3 Apply \\
CO5 & Create production Rust services for web, WebAssembly, embedded or AI targets. & L6 Create \\
\end{cotable}

\capstone{\textbf{Secure High-Performance Inference Service:} a Rust platform with an Axum API, Tokio-based request batching,
an ONNX/Candle model runtime, a PyO3 module for data scientists and a WebAssembly client, delivered with fuzz tests and supply-chain checks.}

\csec{Aligned Careers}
{\raggedright
\role{Rust Systems Engineer}\role{Backend / Infrastructure Engineer}\role{Embedded Software Engineer}
\role{Security Engineer}\role{WebAssembly Developer}\role{AI Infrastructure Engineer}
\par}

\Needspace{12\baselineskip}\csec{Where Rust Fits: 2026--2035 Outlook}
\begin{outlooktable}
Memory-safe systems & Security guidance increasingly favours memory-safe languages for new critical code; Rust is the leading choice. \\
AI infrastructure & Tokenisers, inference engines, data engines (Polars) and vector databases. \\
Edge, embedded \& automotive & Safe firmware and safety-critical components. \\
WebAssembly \& blockchain & Browser, serverless and smart-contract platforms. \\
\end{outlooktable}

\Needspace{15\baselineskip}\csec{Recommended Industry Certifications}
\begin{certtable}
LFEL1002: Getting Started with Rust (exam-based badge) & Linux Foundation & Foundational & I, II \\
LFD480: Programming in Rust (course certificate and badge) & Linux Foundation & Advanced & I--V \\
Courses from Rust Foundation Trusted Training providers & Rust Foundation accredited providers & Course certificate & I--V \\
Ardan Labs Rust certification (vendor exam) & Ardan Labs & Professional & II--IV \\
\end{certtable}
{\footnotesize\itshape Note: there is no official Rust Foundation certification exam yet; in June 2026 the Rust Foundation launched a Trusted Training accreditation for training providers.}

\csec{Text \& Reference Books}
\begin{refs}
\item S. Klabnik and C. Nichols, \emph{The Rust Programming Language}, 2nd ed., No Starch Press, 2023.
\item J. Blandy, J. Orendorff and L. F. S. Tindall, \emph{Programming Rust}, 2nd ed., O'Reilly, 2021.
\item J. Gjengset, \emph{Rust for Rustaceans}, No Starch Press, 2021.
\item L. Palmieri, \emph{Zero To Production in Rust}, 2022.
\end{refs}

%% ============================ COURSE 4: SCALA ========================
\course{4}{EPL-204}{Scala Programming for Big Data \& Functional Systems}{3--0--2--4}{Java or Python, SQL basics}
{Functional programming on the JVM for data platforms, streaming and reliable distributed services}

\csec{Core Competencies}
\begin{itemize}
\item Write expressive Scala 3 combining object-oriented and functional programming.
\item Model domains safely with immutable data, algebraic data types, pattern matching and type classes.
\item Build large-scale batch and streaming data pipelines with Apache Spark, Kafka and Flink.
\item Build reliable concurrent services with effect systems (Cats Effect, ZIO) and actors (Apache Pekko).
\end{itemize}

\csec{Five-Unit Curriculum}
\unit{I}{Scala 3 Foundations on the JVM}{9}{%
Scala and the JVM; sbt and scala-cli; expressions, values and types; methods and functions; classes, objects and traits;
case classes and companion objects; Scala 3 syntax (indentation, enums, extension methods); string interpolation;
interoperability with Java libraries.}
{Builds a student-records application with case classes, enums and traits, using a Java library from Scala.}
\unit{II}{Functional Programming in Scala}{9}{%
Immutability and pure functions; higher-order functions and currying; algebraic data types and pattern matching;
Option, Either and Try for safe error handling; recursion and tail calls; type classes with givens and using clauses;
for-comprehensions; property-based testing with ScalaCheck.}
{Delivers a pure-functional order-validation module with ADTs, Either-based errors and ScalaCheck property tests.}
\unit{III}{Collections, Effects and Concurrency}{9}{%
Immutable and mutable collections; lazy views; Futures and Promises; effect systems: Cats Effect IO and ZIO;
resource safety, fibers and structured concurrency; streaming with fs2 / ZIO Streams; the actor model with Apache Pekko (Akka fork).}
{Builds a concurrent file-processing pipeline in Cats Effect with fs2 streaming and safe resource handling.}
\unit{IV}{Big Data with Apache Spark}{9}{%
Spark architecture (driver, executors, DAG); RDDs, DataFrames and Datasets in Scala; Spark SQL and Catalyst optimiser;
partitioning, joins and performance tuning; Delta Lake / Iceberg tables; Spark Structured Streaming with Kafka;
Spark MLlib pipelines.}
{Ships a Spark-Scala lakehouse job that processes millions of transactions, plus a Structured Streaming fraud-alert pipeline from Kafka.}
\unit{V}{Production Scala: Services, Streaming and Data Platforms}{9}{%
HTTP services with http4s, ZIO HTTP or Play; JSON with Circe; database access with Doobie/Skunk; Kafka producers and consumers;
Apache Flink stream processing concepts; testing (MUnit, ScalaTest, Testcontainers); packaging, Docker and deployment;
observability and JVM tuning.}
{Deploys a containerised http4s service on Kafka and PostgreSQL with full tests and a Spark job orchestrated on a schedule.}

\Needspace{16\baselineskip}\csec{Industry Tools \& Ecosystem}
\begin{toolstable}
Toolchain & Scala 3, sbt, scala-cli, Metals, IntelliJ IDEA, Scalafmt, Scalafix \\
Functional \& effects & Cats, Cats Effect, ZIO, fs2, Circe, ScalaCheck \\
Big data & Apache Spark, Delta Lake, Apache Iceberg, Spark MLlib, Databricks \\
Streaming \& messaging & Apache Kafka, Apache Flink, Apache Pekko (Akka fork) \\
Services \& data access & http4s, ZIO HTTP, Play Framework, Doobie, Skunk \\
Testing \& delivery & MUnit, ScalaTest, Testcontainers, Docker, GitHub Actions \\
\end{toolstable}

\csec{Hands-on Laboratory (30 Hours)}
\labs
{sbt projects, case classes and traits}
{Pattern matching and ADT modelling}
{Type classes, givens and for-comprehensions}
{Cats Effect / ZIO concurrent programs}
{Spark DataFrames and Spark SQL in Scala}
{Spark Structured Streaming with Kafka}
{http4s REST service with Circe and Doobie}
{Containerised Scala service with Testcontainers}

\Needspace{12\baselineskip}\csec{Course Outcomes}
\begin{cotable}
CO1 & Explain Scala 3's type system, syntax and JVM interoperability. & L2 Understand \\
CO2 & Apply functional programming to model domains safely. & L3 Apply \\
CO3 & Build concurrent, resource-safe programs with effect systems and actors. & L3 Apply \\
CO4 & Analyse and optimise large-scale Spark batch and streaming jobs. & L4 Analyse \\
CO5 & Create production Scala data platforms and services. & L6 Create \\
\end{cotable}

\capstone{\textbf{Real-Time Financial Data Platform:} Kafka ingestion, Spark Structured Streaming for fraud detection, Delta Lake storage,
and a functional http4s API serving risk scores, built with typed, tested Scala and deployed in containers.}

\csec{Aligned Careers}
{\raggedright
\role{Big Data Engineer}\role{Scala Developer}\role{Spark Engineer}\role{Streaming Engineer}
\role{Functional Programmer (FinTech)}\role{Data Platform Engineer}
\par}

\Needspace{12\baselineskip}\csec{Where Scala Fits: 2026--2035 Outlook}
\begin{outlooktable}
Big data engines & Spark, Kafka-based systems and many data platforms are built on the JVM; Scala skills help with core engine work. \\
Streaming \& FinTech & Type-safe, functional code for real-time, high-reliability financial systems. \\
Functional programming & Its ideas (immutability, ADTs, effects) shape modern Java, Kotlin, Rust and TypeScript. \\
Data for AI & Large-scale feature pipelines and data preparation for AI products. \\
\end{outlooktable}

\Needspace{15\baselineskip}\csec{Recommended Industry Certifications}
\begin{certtable}
Databricks Certified Associate Developer for Apache Spark & Databricks & Associate & IV \\
Databricks Certified Data Engineer Associate / Professional & Databricks & Assoc./Prof. & IV, V \\
Confluent Certified Developer for Apache Kafka & Confluent & Associate & IV, V \\
Functional Programming in Scala Specialization & EPFL (Coursera) & Course certificate & I--III \\
AWS Certified Data Engineer -- Associate & AWS & Associate & IV, V \\
\end{certtable}

\csec{Text \& Reference Books}
\begin{refs}
\item M. Odersky, L. Spoon, B. Venners and F. Sommers, \emph{Programming in Scala}, 5th ed., Artima, 2021.
\item M. Pilquist, R. Bjarnason and P. Chiusano, \emph{Functional Programming in Scala}, 2nd ed., Manning, 2023.
\item B. Chambers and M. Zaharia, \emph{Spark: The Definitive Guide}, O'Reilly, 2018.
\item D. Wampler, \emph{Programming Scala}, 3rd ed., O'Reilly, 2021.
\end{refs}

%% ============================ COURSE 5: JAVA =========================
\course{5}{EPL-205}{Java Programming for Enterprise \& AI Applications}{3--0--2--4}{Basic programming logic, computer fundamentals}
{Modern Java, Spring Boot microservices and AI-enabled enterprise systems}

\csec{Core Competencies}
\begin{itemize}
\item Write modern Java (LTS 21/25) using records, sealed types, pattern matching, generics, lambdas and streams.
\item Build scalable concurrent applications with virtual threads, structured concurrency and JVM tuning.
\item Develop secure, tested Spring Boot microservices with persistence, messaging and observability.
\item Add AI to enterprise Java with Spring AI and LangChain4j, and deploy cloud-native with containers and GraalVM.
\end{itemize}

\csec{Five-Unit Curriculum}
\unit{I}{Java Foundations and Object-Oriented Programming}{9}{%
JDK, JVM and JRE; compilation and bytecode; data types, operators and control flow; arrays and strings;
classes, objects, constructors, encapsulation, inheritance, polymorphism and abstraction; interfaces;
packages and the module system; exception handling; build tools (Maven, Gradle).}
{Builds a console banking system with full OOP design, custom exceptions and a Maven build.}
\unit{II}{Modern Java: Generics, Collections, Lambdas and Streams}{9}{%
Generics and bounded types; the Collections Framework (List, Set, Map, Queue); equals, hashCode and Comparable;
lambdas, functional interfaces and method references; Stream API and collectors; Optional;
records, sealed classes and pattern matching for switch; text blocks; unit testing with JUnit 5 and Mockito.}
{Delivers an employee-analytics module using records, sealed types and the Stream API, with JUnit 5 tests.}
\unit{III}{Concurrency, the JVM and Performance}{9}{%
Threads, executors and CompletableFuture; virtual threads (Project Loom) and structured concurrency;
synchronisation, locks and concurrent collections; memory model; JVM internals: class loading, JIT, garbage collectors (G1, ZGC);
profiling with JDK Flight Recorder; I/O and NIO.}
{Builds a high-throughput order processor comparing platform and virtual threads, profiled with JFR and tuned GC.}
\unit{IV}{Enterprise Java with Spring Boot}{9}{%
Dependency injection and Spring Core; Spring Boot 3; REST APIs with Spring Web; validation; data access with Spring Data JPA and Hibernate;
transactions; Spring Security with OAuth2/JWT; testing with Spring Boot Test and Testcontainers; API documentation (OpenAPI).}
{Ships a secured Spring Boot e-commerce API with JPA, PostgreSQL, OAuth2/JWT and Testcontainers integration tests.}
\unit{V}{Cloud-Native and AI-Enabled Java}{9}{%
Microservices patterns; messaging with Apache Kafka; caching with Redis; observability with Micrometer and OpenTelemetry;
Docker, Kubernetes and GraalVM native images; AI in Java: Spring AI and LangChain4j for LLM calls, RAG and tool use;
security and dependency scanning (OWASP Dependency-Check).}
{Deploys Kafka-connected Spring Boot microservices on Kubernetes with a GraalVM native image and a Spring AI RAG assistant.}

\Needspace{16\baselineskip}\csec{Industry Tools \& Ecosystem}
\begin{toolstable}
Toolchain & JDK 21/25 (LTS), Maven, Gradle, IntelliJ IDEA, JUnit 5, Mockito \\
Enterprise frameworks & Spring Boot 3, Spring Data JPA, Hibernate, Spring Security, Quarkus, Micronaut \\
Messaging \& data & Apache Kafka, Redis, PostgreSQL, Flyway, Testcontainers \\
Cloud-native & Docker, Kubernetes, GraalVM Native Image, Spring Cloud \\
AI in Java & Spring AI, LangChain4j, Deep Java Library (DJL), ONNX Runtime Java \\
Observability \& security & Micrometer, OpenTelemetry, JDK Flight Recorder, OWASP Dependency-Check, SonarQube \\
\end{toolstable}

\csec{Hands-on Laboratory (30 Hours)}
\labs
{OOP design: classes, inheritance and interfaces}
{Exceptions, files and Maven project build}
{Generics, collections and Comparable}
{Lambdas, streams, records and pattern matching}
{Virtual threads vs platform threads benchmark}
{Spring Boot REST API with JPA and PostgreSQL}
{Spring Security (OAuth2/JWT) and Testcontainers}
{Kafka microservices, GraalVM native image and Spring AI}

\Needspace{12\baselineskip}\csec{Course Outcomes}
\begin{cotable}
CO1 & Explain the JVM, object-oriented principles and Java's type system. & L2 Understand \\
CO2 & Apply generics, collections, lambdas and streams in modern Java. & L3 Apply \\
CO3 & Analyse concurrency and JVM performance and tune applications. & L4 Analyse \\
CO4 & Build secure, tested Spring Boot services with persistence. & L3 Apply \\
CO5 & Create cloud-native, AI-enabled Java microservices. & L6 Create \\
\end{cotable}

\capstone{\textbf{AI-Enabled Enterprise Commerce Platform:} Spring Boot microservices for catalogue, orders and payments connected by Kafka,
secured with OAuth2, deployed on Kubernetes with native images, and a Spring AI assistant that answers customer and staff questions from company data.}

\csec{Aligned Careers}
{\raggedright
\role{Java Developer}\role{Spring Boot Microservices Engineer}\role{Full-Stack Java Developer}
\role{Enterprise Solutions Architect}\role{Android / JVM Developer}\role{AI-Enabled Backend Engineer}
\par}

\Needspace{12\baselineskip}\csec{Where Java Fits: 2026--2035 Outlook}
\begin{outlooktable}
Enterprise systems & Banking, insurance, telecom, e-commerce and government systems run on Java; expected to stay the enterprise backbone. \\
Modern Java & Twice-yearly releases and LTS versions (21, 25) keep adding productivity and performance features. \\
AI in the enterprise & Spring AI and LangChain4j bring LLMs and agents into existing Java systems. \\
Cloud \& big data & Native images, Kubernetes and JVM data engines (Kafka, Flink, Spark). \\
\end{outlooktable}

\Needspace{15\baselineskip}\csec{Recommended Industry Certifications}
\begin{certtable}
Oracle Certified Professional: Java SE 21 Developer (1Z0-830) & Oracle & Professional & I--III \\
Oracle Certified Professional: Java SE 17 Developer (1Z0-829) & Oracle & Professional & I--III \\
Spring Certified Professional 2024 (exam 2V0-72.22) & Broadcom (VMware Tanzu) & Professional & IV \\
Certified Kubernetes Application Developer (CKAD) & CNCF / Linux Foundation & Professional & V \\
NPTEL/SWAYAM: Programming in Java & IIT faculty, Govt. of India & Academic & I, II \\
\end{certtable}

\csec{Text \& Reference Books}
\begin{refs}
\item J. Bloch, \emph{Effective Java}, 3rd ed., Addison-Wesley, 2018.
\item C. S. Horstmann, \emph{Core Java, Vol. I and II}, 13th ed., Pearson, 2024.
\item B. Goetz et al., \emph{Java Concurrency in Practice}, Addison-Wesley, 2006.
\item C. Walls, \emph{Spring in Action}, 6th ed., Manning, 2022.
\item J. Boyarsky and S. Selikoff, \emph{OCP Oracle Certified Professional Java SE 21 Developer Study Guide}, Sybex, 2024.
\end{refs}
